\chapter{Project 5: Worlds Worst Piano}

\section{Overview}
Now that you've had some experience with buttons and an OLED screen, let's add a little more complexity
with a few new types of components. In this project you will:
\begin{itemize}
    \item Learn to use a rotary encoder to provide input to your microcontroller
    \item Write code that makes use of hardware interrupts to coordinate the calling of functions
    \item Make beatuiful music with your microcontroller (OK, maybe not beautiful, but at least audible)
\end{itemize}
At the end of this project, your microcontroller should run a MicroPython program that allows you to write a simple
song and play your creation back. Let's get started!
\begin{figure}[H]
\centering
    \includegraphics[width=.6\linewidth]{project_5/success!.jpg}
    \caption{The end result should look something like this}
\end{figure}

\pagebreak

\section{Directions}

\subsection{Creating the circuit}
Using jumper cables, you will be assembling a circuit between your microcontroller, your breadboard, a potentiometer,
a speaker, and OLED screen included in your kit.

\subsubsection{Remove previous components}
Before beginning, remove all components from prior chapters. You may leave the microcontroller attached to the breadboard.

\subsubsection{Attach the microcontroller to the breadboard}
If it's not already, carefully insert the pins at the bottom of your microcontroller into the breadboard. Refer back
to \ref{pinout} for pin labels. When placing the board into the breadboard, make sure that the microcontroller is oriented such that:
\begin{itemize}
    \item The pin labeled \textbf{5V} is inserted in hole at \textbf{Column H, Row 1} of the breadboard (or \textbf{H1}, for short)
    \item The pin labeled \textbf{GPIO2} is inserted in hole \textbf{D1} of the breadboard
    \item The pin labeled \textbf{GPIO20} is inserted in hole \textbf{H7} of the breadboard
    \item the pin labeled \textbf{GPIO21} is inserted in hole \textbf{D7} of the breadboard
\end{itemize}
You may need to apply more pressure than expected to seat the microcontroller properly in the breadboard. When its over, it should look like this:
\begin{figure}[H]
    \centering
    \includegraphics[width=.6\linewidth]{common/microcontroller_seated_in_breadboard.jpg}
    \caption{So far, so good!}
\end{figure}

\subsubsection{Connect the buttons, resistors, and the screen}
\begin{itemize}
    \item Place the OLED screen into the breadboard such that its 4 pins connect to holes \textbf{A11}, \textbf{A12}, \textbf{A13}, and \textbf{A14}.
    \item Place the potentiometer so that it is in \textbf{A25}, \textbf{A26}, \textbf{A27}, \textbf{A28}, and \textbf{A29} and the knob is facing you.
    \item Place a button so that one connected set of pins (refer to \ref{button_basics} for an example) is in \textbf{F27}
    and \textbf{I27} and the other set is in \textbf{F29} and \textbf{I29}.
    \item Finally, insert the speaker wires so that the red is in \textbf{J17} and the black in in \textbf{J18}.
\end{itemize}

You should be left with something that looks like this:
\begin{figure}[H]
    \centering
    \includegraphics[width=.6\linewidth]{project_5/components_placed.jpg}
    \caption{All of the components except for the jumper wires are now placed.}
\end{figure}

\subsubsection{Connect the necessary jumper wires}
\begin{itemize}
    \item Using a black jumper wire, place one end of the wire into hole \textbf{I2} of the breadboard and the other
    end into \textbf{D11}. This will provide the ground connection for the OLED screen.
    \item Place another black jumper wire between \textbf{J2} and \textbf{E25}. This will provide the ground connection
    for the potentiometer.
    \item Place one end of a red jumper wire into hole \textbf{I3} of the breadboard and the other end into
    \textbf{D12}. This will provide \textbf{3.3} volts of power to the OLED screen.
    \item Place another red jumper wire between \textbf{J3} and \textbf{E26}. This will provide the power connection
    for the potentiometer.
    \item Place a white jumper wire between \textbf{J4} and \textbf{E29}. This will provide a clock signal for the potentiometer.
    \item Place a purple jumper wire between \textbf{J5} and \textbf{E28}. This will provide the delta time signal for the potentiometer.
    \item Place a green jumper wire between \textbf{J6} and \textbf{E27}. This will provide the button signal on the potentiometer.
    \item Place a yellow jumper wire between \textbf{J7} and \textbf{G17}. This will power the speaker.
    \item Place a brown jumper wire between \textbf{G18} and \textbf{E11}. This will provide the ground connection for the speaker.
    \item Place a green jumper wire between \textbf{J27} and \textbf{B4}. This will provide the signal for the play button.
    \item Place a brown jumper wire between \textbf{J29} and \textbf{F18}. This will provide ground for the play button.
    \item Place a yellow jumper wire between \textbf{D14} and \textbf{B5}. This will provide the data for the OLED.
    \item Finally, place a white jumper wire between \textbf{D13} and \textbf{B6}. This will provide the clock signal for the OLED.
\end{itemize}

You should be left with something that looks like this:
\begin{figure}[H]
    \centering
    \includegraphics[width=.6\linewidth]{project_5/all_wired_up.jpg}
    \caption{All components and wires are now installed}
\end{figure}

\subsection{Programming the microcontroller}
Once all of the wiring is correct, connect the USB cable to the microcontroller and load the IDE to
access it. Refer back to Chapter \ref{ide} for instructions.

Click on the file named "project\_5\_worlds\_worst\_game.py". This will load the code in the editor for this section. Read through the comments
and the code to get a sense for how it works. Once you are ready, you can click the blue play button in the upper left of the window
to start the game.

The OLED will display the current note selected and turning the rotary encoder will change the note. Pressing the button on the rotary encoder
will add the currently displayed note to your song. Pressing the play button will play back your song through the speaker.

\section{Review}
In this project, we used a rotary encoder to provide a way to select notes and add them to a song. We also used a button to trigger
the playback of the song. The OLED screen provided feedback to the user about which note was currently selected.
Finally, we used a speaker to play back the song.

\section{Possible Extensions}
If you want to do some experimentation, try these:

\begin{itemize}
    \item Provide a fun visualization on the OLED screen when the song is playing.
    \item Add another button that lets you choose between different instrument sounds for each note (hint, change the duty cycle).
\end{itemize}
